\chapter{Acceleration}

By constructing a specifically difficult function to optimize, one can show that for any first order method which satisfies
\begin{equation}
    x_k\in x_0+\mathrm{span}\{\nabla f(x_0),\dots,\nabla f(x_{k-1})\}
\end{equation}
there exists $L$-smooth $f$ such that
\begin{equation}
    f(x_k)-\min f\geq\frac{3L\|x_0-x^*\|}{32(k+1)^2}.
\end{equation}
Note that up until now we have not reached convergence better than $\Oc(1/k)$. Closing the gap to $\Oc(1/k^2)$ is the goal of this section.

\section{Polyak's heavy ball method}

Note that we can interpret gradient descent for minimizing the (unconstrained) problem
\begin{equation}
    \min_x f(x)
\end{equation}
as a discretization of the \emph{gradient flow}
\begin{equation}
    \dot x  = -\nabla f(x).
\end{equation}
Interpreting the function $f$ as a potential, $\nabla f$ corresponds to the force enacted by this potential on a particle (with mass one). It is now natural to include friction into this model. The force caused by friction is typically modelled proportional to the velocity leading to
\begin{equation}\label{acceleration:eq:friction_ode}
    \ddot x = - \gamma \dot x -\nabla f(x).
\end{equation}
Let us discretize this ODE via $\dot x (t)\approx \frac{x(t)-x(t-h)}{h}$ and 
\begin{equation}
    \ddot x(t)
    \approx \frac{\dot x(t+h)-\dot x(t)}{h}
    \approx \frac{x(t+h)-2x(t)+x(t-h)}{h^2}.
\end{equation}
Denoting $x_{t-h}=x_{k-1}$, $x(t)=x_k$ and $x(t+h)=x_{k+1}$ we find via inserting into~\cref{acceleration:eq:friction_ode}
\begin{equation}
    x_{k+1} - 2x_k + x_{k-1} = -h\gamma (x_{k}-x_{k-1}) - h^2\nabla f(x_k)
\end{equation}
which yields after rearranging the update
\begin{equation}
    x_{k+1} = x_{k} + (1-\gamma h)(x_{k}-x_{k-1}) - h^2\nabla f(x_k).
\end{equation}
Relabelling the constants and allowing them also to be iteration depdendent we arrive at \emph{Polyak's heavy ball method}
\begin{equation}
    x_{k+1} = x_{k} + \beta_k(x_{k}-x_{k-1}) - \tau_k\nabla f(x_k).
\end{equation}
Note that even without this derivation using the ODE perspective, the update rule is rather intuitive: The update contains in addition to the descent direction for $f$ also an additional \emph{inertia} term, which adds a compontent in the same direction as the previous step.

The convergence proof of the heavy ball method will rely on the following basic results from linear algebra which we include here for the sake of completeness.

\begin{lemma}\label{acceleration:lemma:spectral_radius}
    For any $A\in \R^{d\times d}$ we have $\lim_{n\rightarrow\infty}\|A^n\|^{1/n} = \rho(A)$ where 
    \begin{equation}
        \rho(A) = \max_{i}|\lambda_i(A)|
    \end{equation}
    with $\lambda_i(A)$ the eigenvalues of $A$.
\end{lemma}
\begin{proof}
    First of all it is easy to see that $\liminf_{n\rightarrow\infty}\|A^n\|^{1/n} \geq \rho(A)$. Indeed, let $v$ be a normalized eigenvector for the largest (in absolute value) eigenvalue of $A$, then we have
    \begin{equation}
        \|A^n\|^{1/n} \geq \|A^n v\|^{1/n} = \rho(A)
    \end{equation}
    and thus also $\liminf_{n\rightarrow\infty}\|A^n\|^{1/n} \geq \rho(A)$.
    The converse inequality is a little more subtle.
    Recall from linar algebra that for every square matrix we can derive the Jordan normal form
    \begin{equation}
        A = B^{-1}D B
    \end{equation}
    with 
    \begin{equation}
        D = \begin{pmatrix}
            \lambda_1& \rho_1 & \dots&\dots&\\
            0 & \lambda_2& \rho_2&\dots&\\
            &&\ddots&\ddots&\\
            & &\dots & \lambda_{d-1}&\rho_{d-1}&\\
            & & &\dots&\lambda_d\\
        \end{pmatrix}
    \end{equation}
    with $\lambda_i$ the eigenvalues (with multiplicity) and $\rho_i\in \{0,1\}$ and which we write as $D=\Lambda + J$ where $\Lambda$ contains the diagonal entries and $J$ the off diagonals. Note that $J^k=0$ whenever $k$ is larger than the largest Jordan block. In particular, $J^k=0$ for $k\geq d-1$. Since $\Lambda$ and $J$ commute, we, thus, find
    \begin{equation}
        \begin{aligned}
            \|D^n\| = 
            \|\sum_{k=0}^{d-2}
            \begin{pmatrix}
                n\\
                k
            \end{pmatrix}\Lambda^{n-k}J^k\|
            \leq \sum_{k=0}^{d-2}
            \begin{pmatrix}
                n\\
                k
            \end{pmatrix}\|\Lambda^{n-k}\|
            \leq& \sum_{k=0}^{d-2} \frac{n^k}{k!}\|\Lambda^{n-k}\|\\
            \leq& \sum_{k=0}^{d-2} \frac{n^k}{k!}\rho(A)^{n-k}\\
            =& \rho(A)^{n}\sum_{k=0}^{d-2} \frac{n^k}{k!}\rho(A)^{-k}\\
        \end{aligned}
    \end{equation}
    which implies
    \begin{equation}
        \begin{aligned}
            \|D^n\|^{1/n}
            \leq& \rho(A)\left(\sum_{k=0}^{d-2} \frac{n^k}{k!}\rho(A)^{-k}\right)^{1/n}
            \leq& \rho(A)\left((d-1)n^{d-2}\max_{k=0,\dots, d-2}\rho(A)^{-k}\right)^{1/n}
        \end{aligned}
    \end{equation}
    where we note that tha maximizing $k$ above is either $k=0$ or $k=d-2$ when $\rho(A)\leq 1$ or $\rho(A)>1$, respectively. Since $\lim_{n\rightarrow \infty} n^{1/n}=1$ it follows \begin{equation}
        \begin{aligned}
            \limsup_n\|D^n\|^{1/n}
            \leq \rho(A).
        \end{aligned}
    \end{equation}
    Lastly, we conclude
    \begin{equation}
        \limsup_n\|A^n\|^{1/n} 
        = \limsup_n \|B^{-1}D^nB\|^{1/n}
        \leq \limsup_n\|B^{-1}\|^{1/n}\|D^n\|^{1/n} \|B\|^{1/n}\leq \rho(A)
    \end{equation}
\end{proof}

\begin{cor}\label{acceleration:cor:spectral_radius}
    It holds that $\lim_k A^k=0$ if and only if $\rho(A)<1$. Moreover, in this case for every $\epsilon>0$ there exists $c(\epsilon)>0$ such that $\|A^k\|\leq c(\epsilon)(\rho(A)+\epsilon)^k$.
\end{cor}
\begin{proof}
    The \enquote{if and only if}-part of the theorem follows from directly from~\cref{acceleration:lemma:spectral_radius}.
    Moreover, by~\cref{acceleration:lemma:spectral_radius} we can choose $n$ sufficiently large such that $\|A^k\|^{1/k}\leq \rho(A) + \epsilon$ for $k\geq n$.
    For such $k$ it follows
    \begin{equation}
        \|A^k\|\leq (\rho(A)+\epsilon)^k.
    \end{equation}
    Since there also exists some $c>0$ such that 
    \begin{equation}
        \|A^k\|\leq c(\rho(A)+\epsilon)^k.
    \end{equation}
    for $k\leq n$ the result follows.
\end{proof}

We can now prove convergence of the heavy ball method. The technique we use is quite standard: We analyse the eigenvalues of the linearization of the update rule.

\begin{theorem}[Convergence of heavy ball]
    Let 
    \begin{equation}
        0\leq \beta <1, \quad 0<\tau <2(1+\beta)/L,\quad \mu \Id\preceq \nabla^2 f(x^*)\preceq L\Id.
    \end{equation}
    There exists $\epsilon>0$ such that for $x_0,x_1\in B_\epsilon(x^*)$ it holds that $x_k\rightarrow x^*$. More specifically, 
    \begin{equation}
        \|x_k-x^*\|\leq c(\delta)(q+\delta)^k,\quad 0\leq q< 1, 0<\delta<1-q.
    \end{equation}
    Moreover, we can derive the optimal parameters (\ie, the minimal value of $q$) as
    \begin{equation}
        q = \frac{\sqrt{L}-\sqrt{\mu}}{\sqrt{L}+\sqrt{\mu}},
        \quad \tau = \frac{4}{(\sqrt{L} + \sqrt{\mu})^2}
        ,\quad \beta = \left(\frac{\sqrt{L}-\sqrt{\mu}}{\sqrt{L}+\sqrt{\mu}}\right)^2
    \end{equation}
\end{theorem}
\begin{proof}
    Note that, as is, the scheme is, in fact, a recursion of depth two. We, thus, rewrite the update as
    \begin{equation}
        \begin{aligned}
            \begin{bmatrix}
                x_{k+1}\\
                x_k
            \end{bmatrix}
            &= 
            \begin{bmatrix}
                x_{k} + \beta(x_{k}-x_{k-1}) - \tau\nabla f(x_k)\\
                x_k
            \end{bmatrix}\\
            &=
            \begin{pmatrix}
                1+\beta - \tau \nabla f&-\beta\\
                1&0
            \end{pmatrix}
            \begin{bmatrix}
                x_k\\
                x_{k-1}
            \end{bmatrix}.
        \end{aligned}
    \end{equation}
    Moreover, we can subtract the minimizer $x^*$, \ie, 
    \begin{equation}
        \begin{aligned}
            \begin{bmatrix}
                x_{k+1}-x^*\\
                x_k-x^*
            \end{bmatrix}
            &=
            \begin{pmatrix}
                1+\beta - \tau \nabla f&-\beta\\
                1&0
            \end{pmatrix}
            \begin{bmatrix}
                x_k-x^*\\
                x_{k-1}-x^*
            \end{bmatrix}.
        \end{aligned}
    \end{equation}
    Moreover, we can \emph{linearize} the update as follows: By the fundamental theorem of calculus we have
    \begin{equation}
        \begin{aligned}
            \nabla f(x_k) - \nabla f(x^*) 
            =& \int_0^1 \nabla^2 f(x^* + t(x_k-x^*))\dd t (x_k-x^*)\\
            =& \nabla^2 f(x^*)(x_k-x^*) + \int_0^1 \nabla^2 f(x^* + t(x_k-x^*))-\nabla^2 f(x^*)\dd t (x_k-x^*).
        \end{aligned}
    \end{equation}
    Thus, defining $z_k = \int_0^1 \nabla^2 f(x^* + t(x_k-x^*))-\nabla^2 f(x^*)\dd t (x_k-x^*)$ we have
    \begin{equation}
        \begin{aligned}
            \begin{bmatrix}
                x_{k+1}-x^*\\
                x_k-x^*
            \end{bmatrix}
            &=
            \underbrace{\begin{pmatrix}
                1+\beta - \tau \nabla^2 f(x^*)&-\beta\\
                1&0
            \end{pmatrix}}_{\mathcal{H}}
            \begin{bmatrix}
                x_k-x^*\\
                x_{k-1}-x^*
            \end{bmatrix}
            + \begin{bmatrix}
                z_k\\
                0
            \end{bmatrix}
        \end{aligned}
    \end{equation}
    In view of~\cref{acceleration:cor:spectral_radius}, we analyze the eigenvalues of $\mathcal{H}$. Let $(v,w)$ be an eigenvector of $\mathcal{H}$ with eigenvalue $\lambda$. We find
    \begin{equation}
        \begin{cases}
            (1+\beta -\tau \nabla^2 f(x^*))v - \beta w &= \lambda v\\
            v &= \lambda w.
        \end{cases}
    \end{equation}
    Inserting the second into the first equation it follows
    \begin{equation}
        (1+\beta -\tau \nabla^2 f(x^*))v = \lambda v + \beta\lambda^{-1} v.
    \end{equation}
    This, in turn, implies that there exist an eigenvalue $\eta$ of $\nabla^2f(x^*)$ such that
    \begin{equation}
        1+\beta -\tau \eta = \lambda + \beta\lambda^{-1},
    \end{equation}
    respectively
    \begin{equation}
        \lambda^2 - \lambda(1+\beta -\tau \eta) = -\beta
    \end{equation}
    by completing the square we obtain
    \begin{equation}\label{acceleration:eq:heavy_ball1}
        \left(\lambda - \frac{1+\beta -\tau \eta}{2}\right)^2 = \left(\frac{1+\beta -\tau \eta}{2}\right)^2 - \beta.
    \end{equation}
    Our goal is to ensure $|\lambda|<1$. Note that~\eqref{acceleration:eq:heavy_ball1} shows that $\lambda$ is contained within a circle in the complex plane with center $c \coloneqq \frac{1+\beta -\tau \eta}{2}$ and radius $\sqrt{c^2 - \beta}$. It is easy to see that it suffices to estimate the values of $\lambda$ for maximal and minimal $c$.
    A strict upper bound for $c$ is obtained by the estimate $\tau\eta>0$ yielding $c = (1+\beta)/2$ and 
    \begin{equation}
        \begin{aligned}
            |\lambda|\leq c+c^2-\beta = \frac{3+\beta^2}{4}<1
        \end{aligned}
    \end{equation}
    due to $0\leq \beta<1$.
    Conversely, a strict lower bound for $c$ is obtained for $\tau = 2(1+\beta)/L$ and $\eta = L$ leading to $c = -(1+\beta)$ and a similar estimate as above. Thus, we find $\rho(\mathcal{H})<1$. Noting that $z_k = o(x_k-x^*)$ the result follows where we leave the remaining details as an exercise.
\end{proof}

% \todo{@Alex: Exercise: prove local convergence from the linearization with hints! Prove optimal parameters with hints.}

\section{Nesterov acceleration}
Polyak's method is in fact optimal in the sense that it achieves the best possible convergence rate. However, this is only true for strongly convex and twice continuously differentiable $f$. The next acceleration we consider achieves optimal convergence under significantly weaker conditions. We consider now again problems of the form 
\begin{equation}
    \min_x F(x)\coloneqq f(x) + g(x)
\end{equation}
where only $f$ is $L$-smooth. The update of the the \enquote{fast iterative shrinkage-thresholding algorithm} (FISTA, or fast proximal gradient method) reads as
\begin{equation}
    \begin{cases}
        t_{k+1} = \frac{1+\sqrt{1+4t_k^2}}{2}\\
        \beta_k = \frac{t_k-1}{t_{k+1}}\\
        y_k = x_k + \beta_k(x_k-x_{k-1})\\
        x_{k+1} = \prox_{\tau g}(y_k - \tau \nabla f(y_k)).
    \end{cases}
\end{equation}
In the case $g\equiv 0$ the method closely resebles the heavy ball algorithm. However, the gradient is now evaluated after adding the inertia term. Moreover the convergence relies on a subtle adaptive choice of the inertia parameter.

\begin{lemma}
    Set $t_1 = 1$. It holds true that 
    \begin{equation}
        t_k \geq \frac{k+1}{2}
    \end{equation}
    for all $k$.
\end{lemma}
\begin{proof}
    The proof follows by straightforward induction. The assertion holds true for $k=1$ and, using the induction hypothesis, we find
    \begin{equation}
        t_{k+1} = \frac{1+\sqrt{1+4t_k^2}}{2}\geq \frac{1+\sqrt{1+(k+1)^2}}{2}\geq \frac{1+(k+1)}{2}.
    \end{equation}
\end{proof}

\begin{lemma}[Fundamental prox-grad inequality]
    Denote 
    \begin{equation}
        P_\tau(x) = \prox_{\tau g}(x-\tau \nabla f(x)).
    \end{equation}
    For any $\tau\leq \frac{1}{L}$ it holds true that 
    \begin{equation}
        F(x) - F(P_\tau (y))\geq \frac{1}{2\tau}\|x-P_\tau (y)\|^2 - \frac{1}{2\tau}\|x-y\|^2 + \ell_f(x,y)
    \end{equation}
    where $\ell_f(x,y) = f(x) - f(y) - \langle x-y,\nabla f(y)\rangle$
\end{lemma}
\begin{proof}
    First note that we can write
    \begin{equation}
        \begin{aligned}
            P_\tau (y) 
            =& \arg\min_z \frac{1}{2\tau}\|z-(y-\tau \nabla f(y))\|^2 + g(z)\\
            =& \arg\min_z f(y) + \langle z-y,\nabla f(y)\rangle + \frac{1}{2\tau}\|z-y\|^2 + g(z).
        \end{aligned}
    \end{equation}
    As a sidenote, this shows that the proximal-gradient method can be interpreted as a proximal method of $g$ plus a linear approximation of $f$. Let us denote 
    \begin{equation}
        \phi(z) = f(y) + \langle z-y,\nabla f(y)\rangle + \frac{1}{2\tau}\|z-y\|^2 + g(z).
    \end{equation}
    By convexity of and $g$, $\phi$ is $1/\tau$-strongly convex and we find
    \begin{equation}
        \phi(x) - \phi(P_\tau (y))\geq \frac{1}{2\tau}\|x-P_\tau (y)\|^2
    \end{equation}
    Moreover, by convexity and $L$-smoothness of $f$ and the fact that $\frac{1}{\tau}\geq L$
    \begin{equation}
        f(z) \leq f(y) + \langle z-y,\nabla f(y)\rangle + \frac{L}{2}\|y-z\| \leq f(y) + \langle z-y,\nabla f(y)\rangle + \frac{1}{2\tau}\|z-y\|^2
    \end{equation}
    and thus
    \begin{equation}
        \phi(x) - F(P_\tau (y))\geq \frac{1}{2\tau}\|x-P_\tau (y)\|^2
    \end{equation}
    Pugging in $\phi$ yields
    \begin{equation}
        f(y) + \langle x-y,\nabla f(y)\rangle + \frac{1}{2\tau}\|x-y\|^2 + g(x) - F(P_\tau (y))\geq \frac{1}{2\tau}\|x-P_\tau (y)\|^2
    \end{equation}
    concluding the proof.
\end{proof}
\begin{theorem}[Convergence of FISTA]
    Let the iterates of FISTA be initialized as $x_0=x_1$ and $t_1=1$ and the step size $\tau\leq \frac{1}{L}$ with $L$ the Lipschitz constant of $\nabla f$. Then we obtain the convergence 
    \begin{equation}
        F(x_k) - F(x^*)\leq \frac{C}{k^2 \tau}
    \end{equation}
    with a constnat $C$ depending on the initialization.
\end{theorem}
\begin{proof}
    Choosing $x = t_{k+1}^{-1}x^* + (1-t_{k+1}^{-1})x_k$ and $y=y_k$ in the fundamental prox-grad inequality and noting that by convexity of $f$ $\ell_f(x,y)\geq 0$, we obtain
    \begin{equation}\label{eq:fista1}
        \begin{aligned}
            F(t_{k+1}^{-1}x^* + &(1-t_{k+1}^{-1})x_k) - F(x_{k+1})\\
            \geq& \frac{1}{2\tau}\|t_{k+1}^{-1}x^* + (1-t_{k+1}^{-1})x_k-x_{k+1}\|^2 - \frac{1}{2\tau}\|t_{k+1}^{-1}x^* + (1-t_{k+1}^{-1})x_k-y_k\|^2\\
            \geq& \frac{1}{2\tau t_{k+1}^2}\|x^* + (t_{k+1}-1)x_k-t_{k+1} x_{k+1}\|^2 - \frac{1}{2\tau t_{k+1}^2}\|x^* + (t_{k+1}-1)x_k-t_{k+1}y_k\|^2.
        \end{aligned}
    \end{equation}
    Again, by convexity and the fact that by induction one easily verifies $t_k\geq 1$ for all $k$, it holds true that
    \begin{equation}\label{eq:fista2}
        \begin{aligned}
            F(t_{k+1}^{-1}x^* + (1-t_{k+1}^{-1})x_k) - F(x_{k+1})
            \leq& t_{k+1}^{-1} F(x^*) + (1-t_{k+1}^{-1}) F(x_k) - F(x_{k+1})\\
            \leq& (1-t_{k+1}^{-1}) (F(x_k)-F(x^*)) - (F(x_{k+1})-F(x^*)).
        \end{aligned}
    \end{equation}
    Since $y_k = x_k + \frac{t_k-1}{t_{k+1}}(x_k-x_{k-1})$ we have
    \begin{equation}\label{eq:fista3}
        \begin{aligned}
            \|x^* + (t_{k+1}-1)x_k-t_{k+1}y_k\|^2
            = \|x^* - t_k x_k + (t_k-1)x_{k-1}\|^2
        \end{aligned}
    \end{equation}
    Combining~\eqref{eq:fista1}, \eqref{eq:fista2}, and \eqref{eq:fista3} yields
    \begin{equation}
        \begin{aligned}
            t_{k+1}^2(1-t_{k+1}^{-1}) &(F(x_k)-F(x^*)) - t_{k+1}^2(F(x_{k+1})-F(x^*))\\
            \geq& \frac{1}{2\tau}\|x^* - t_{k+1} x_{k+1}+ (t_{k+1}-1)x_k\|^2 - \frac{1}{2\tau}\|x^* - t_k x_k + (t_k-1)x_{k-1}\|^2.
        \end{aligned}
    \end{equation}
    By the update rule of FISTA, it holds true that $t_{k+1}^2(1-t_{k+1}^{-1}) = t_{k}^2$ and we have
    \begin{equation}
        \begin{aligned}
            t_{k}^2 &(F(x_k)-F(x^*)) - t_{k+1}^2(F(x_{k+1})-F(x^*))\\
            \geq& \frac{1}{2\tau}\|x^* - t_{k+1} x_{k+1}+ (t_{k+1}-1)x_k\|^2 - \frac{1}{2\tau}\|x^* - t_k x_k + (t_k-1)x_{k-1}\|^2
        \end{aligned}
    \end{equation}
    hence,
    \begin{equation}
        \begin{aligned}
            t_{k}^2 &(F(x_k)-F(x^*)) + \frac{1}{2\tau}\|x^* - t_k x_k + (t_k-1)x_{k-1}\|^2\\
            \geq& t_{k+1}^2(F(x_{k+1})-F(x^*)) + \frac{1}{2\tau}\|x^* - t_{k+1} x_{k+1}+ (t_{k+1}-1)x_k\|^2
        \end{aligned}
    \end{equation}
    and by iterating over $k$ and noting that $x_0=x_1$
    \begin{equation}
        \begin{aligned}
            t_{2}^2 &(F(x_2)-F(x^*)) + \frac{1}{2\tau}\|t_2 (x^* - x_2) + (t_2-1)(x_{1}-x^*)\|^2\\
            \geq& t_{k}^2(F(x_{k})-F(x^*)) + \frac{1}{2\tau}\|x^* - t_{k} x_{k}+ (t_{k}-1)x_{k-1}\|^2
        \end{aligned}
    \end{equation}
    Moreover, the fundamental prox-grad inequality with $x=x^*$ and $y=y_1=x_0=x_1$ yields
    \begin{equation}
        F(x^*) - F(x_2)\geq \frac{1}{2\tau}\|x^* - x_2\|^2 - \frac{1}{2\tau}\|x^* - x_1\|^2.
    \end{equation}
    Thus, we find
    \begin{equation}
        \begin{aligned}
            F(x_{k})-F(x^*)
            \leq t_k^{-2}\left(\frac{t_2^2}{2\tau}\|x^* - x_1\|^2 + \frac{1}{2\tau}\|t_2 (x^* - x_2) + (t_2-1)(x_{1}-x^*)\|^2\right)
        \end{aligned}
    \end{equation}
    Since $t_k\geq \frac{k+1}{2}$ the result follows.
\end{proof}